\documentclass[12pt, a4paper]{article}
\usepackage{graphicx} % Required for inserting images
\graphicspath{{images/}{images-desafios/}}
\usepackage[a4paper, total={7.5in, 10in}]{geometry}
\usepackage{float} % pra conseguir usar o \begin{figure}[H]
\usepackage{indentfirst} % pra fazer tab n 1 paragrafos dos capitulos
% package pra tabela
\usepackage{multirow}
\usepackage{array}

%pacote pra codigo
\usepackage{listings} 
\usepackage{xcolor}
\definecolor{codegreen}{rgb}{0,0.6,0}
\definecolor{codegray}{rgb}{0.5,0.5,0.5}
\definecolor{codepurple}{rgb}{0.58,0,0.82}
\definecolor{backcolour}{rgb}{0.95,0.95,0.92}

\lstdefinestyle{myCodeStyle}{
	backgroundcolor=\color{backcolour},   
	commentstyle=\color{codegreen},
	keywordstyle=\color{cyan},
	numberstyle=\tiny\color{codegray},
	stringstyle=\color{codepurple},
	basicstyle=\ttfamily\footnotesize,
	breakatwhitespace=false,         
	breaklines=true,                 
	captionpos=b,                    
	keepspaces=true,                 
	%numbers=left,                    
	numbersep=5pt,                  
	showspaces=false,                
	showstringspaces=false,
	showtabs=false,                  
	tabsize=2
}
\lstset{style=myCodeStyle}
\renewcommand{\lstlistingname}{Código}%Muda o termo "Listing":  Listing -> Código
\renewcommand{\lstlistlistingname}{Lista de \lstlistingname s}% List of Listings -> Lista de Códigos


\newcommand{\unit}[1]{\ensuremath{\, \mathrm{#1}}} % comando pra tirar do math mode a unidade de medida

\def\tamanhoGrafico{0.8}

\newcommand{\ohm}{\unit{\Omega} }
\newcommand{\kohm}{\unit{k\Omega}}

\title{
	Laboratório Digital A - PCS3335
	\\ Planejamento da Experiência 8
}

\author{
	Alunos:
	\\ Rafael Hipólito Rodrigues - N.USP: 14604308
	\\ {{PERSON-2}} - N.USP: {{PERSON-ID-2}}
	\\ \\
	Turma 3 - Bancada A6
	\\ Professor: Bruno de Carvalho Albertini
}

%\date{}

\begin{document}
	
	\maketitle
	
	\section*{Planejamento} 
	
	Finalização da UART simplificada. Elaboração da entidade UART de acordo com o datasheet do PC16550D.
	Essa entidada possou todas as entidades de transmissão e recepção desenvolvidas nas experiências anteriores.
	
	\section*{Diagramas}
	
	\subsection*{FSM}
	
	Diagrama da Maquina de Estados Finitos do \texttt{Transmitter} gerado no Quartus : 
	
	\begin{figure}[H]
		\centering
		\includegraphics[width=0.9\linewidth]{images/FSM_transmitter.png}
		\caption{FSM do Transmitter.}
		\label{fig:FSM_transmitter}
	\end{figure}
	
	Diagrama da Maquina de Estados Finitos do \texttt{Receiver} gerado no Quartus : 
	
	\begin{figure}[H]
		\centering
		\includegraphics[width=0.9\linewidth]{images/FSM_receiver.png}
		\caption{FSM do Receiver.}
		\label{fig:FSM_receiver}
	\end{figure}
	
	
	
	\subsection*{RTL}
	
	Diagramas RTL da entidades gerado no Quartus:
	
	\begin{figure}[H]
		\centering
		\includegraphics[width=0.9\linewidth]{images/RTL_topLevel.png}
		\caption{Diagrama RTL do top level.}
		\label{fig:RTL_topLevel}
	\end{figure}
	
	\begin{figure}[H]
		\centering
		\includegraphics[width=0.9\linewidth]{images/RTL_uart.png}
		\caption{Diagrama RTL do \texttt{UART}.}
		\label{fig:RTL_uart}
	\end{figure}

	
	Os contadores e registradores utilizados foram os mesmos da experiência 2.
	
	\section*{Pinagem}

	
	\begin{table}[H]
		\begin{center}
			\begin{tabular}{|l|l|l|l|}
				\hline
				serialIn              & PIN\_N16  & Input  & GPIO 0           \\ \hline
				serialOut             & PIN\_B16  & Output & GPIO 1           \\ \hline
				A{[}0{]}              & PIN\_M16  & Input  & GPIO 2           \\ \hline
				A{[}1{]}              & PIN\_D17  & Input  & GPIO 3           \\ \hline
				A{[}2{]}              & PIN\_K21  & Input  & GPIO 4           \\ \hline
				notADS                & PIN\_M20  & Input  & GPIO 5           \\ \hline
				notBAUDOUT            & PIN\_N21  & Output & GPIO 6           \\ \hline
				notRXRDY              & PIN\_R21  & Output & GPIO 7           \\ \hline
				notTXRDY              & PIN\_N20  & Output & GPIO 8           \\ \hline
				RD                    & PIN\_M22  & Input  & GPIO 9           \\ \hline
				reset                 & PIN\_L22  & Input  & GPIO 10          \\ \hline
				WR                    & PIN\_P16  & Input  & GPIO 11          \\ \hline
				clock50M              & PIN\_M9   & Input  & FPGA clock 50MHz \\ \hline
				leds{[}0{]}           & PIN\_AA2  & Output & leds{[}0{]}      \\ \hline
				leds{[}1{]}           & PIN\_AA1  & Output & leds{[}1{]}      \\ \hline
				leds{[}2{]}           & PIN\_W2   & Output & leds{[}2{]}      \\ \hline
				leds{[}3{]}           & PIN\_Y3   & Output & leds{[}3{]}      \\ \hline
				leds{[}4{]}           & PIN\_N2   & Output & leds{[}4{]}      \\ \hline
				leds{[}5{]}           & PIN\_N1   & Output & leds{[}5{]}      \\ \hline
				leds{[}6{]}           & PIN\_U2   & Output & leds{[}6{]}      \\ \hline
				leds{[}7{]}           & PIN\_U1   & Output & leds{[}7{]}      \\ \hline
				leds{[}8{]}           & PIN\_L2   & Output & leds{[}8{]}      \\ \hline
				leds{[}9{]}           & PIN\_L1   & Output & leds{[}9{]}      \\ \hline
				display7segRBR{[}0{]} & PIN\_AA20 & Output & Display 7 seg 2  \\ \hline
				display7segRBR{[}1{]} & PIN\_AB20 & Output & Display 7 seg 2  \\ \hline
				display7segRBR{[}2{]} & PIN\_AA19 & Output & Display 7 seg 2  \\ \hline
				display7segRBR{[}3{]} & PIN\_AA18 & Output & Display 7 seg 2  \\ \hline
				display7segRBR{[}4{]} & PIN\_AB18 & Output & Display 7 seg 2  \\ \hline
				display7segRBR{[}5{]} & PIN\_AA17 & Output & Display 7 seg 2  \\ \hline
				display7segRBR{[}6{]} & PIN\_U22  & Output & Display 7 seg 2  \\ \hline
				display7segTHR{[}0{]} & PIN\_U21  & Output & Display 7 seg 1  \\ \hline
				display7segTHR{[}1{]} & PIN\_V21  & Output & Display 7 seg 1  \\ \hline
				display7segTHR{[}2{]} & PIN\_W22  & Output & Display 7 seg 1  \\ \hline
				display7segTHR{[}3{]} & PIN\_W21  & Output & Display 7 seg 1  \\ \hline
				display7segTHR{[}4{]} & PIN\_Y22  & Output & Display 7 seg 1  \\ \hline
				display7segTHR{[}5{]} & PIN\_Y21  & Output & Display 7 seg 1  \\ \hline
				display7segTHR{[}6{]} & PIN\_AA22 & Output & Display 7 seg 1  \\ \hline
				switches{[}0{]}       & PIN\_U13  & Input  & switches{[}0{]}  \\ \hline
				switches{[}1{]}       & PIN\_V13  & Input  & switches{[}1{]}  \\ \hline
				switches{[}2{]}       & PIN\_T13  & Input  & switches{[}2{]}  \\ \hline
				switches{[}3{]}       & PIN\_T12  & Input  & switches{[}3{]}  \\ \hline
				switches{[}4{]}       & PIN\_AA15 & Input  & switches{[}4{]}  \\ \hline
				switches{[}5{]}       & PIN\_AB15 & Input  & switches{[}5{]}  \\ \hline
				switches{[}6{]}       & PIN\_AA14 & Input  & switches{[}6{]}  \\ \hline
				switches{[}7{]}       & PIN\_AA13 & Input  & switches{[}7{]}  \\ \hline
				switches{[}8{]}       & PIN\_AB13 & Input  & switches{[}8{]}  \\ \hline
				switches{[}9{]}       & PIN\_AB12 & Input  & switches{[}9{]}  \\ \hline
			\end{tabular}
		\end{center}
	\end{table}
	
	
	
	\section*{Testes}
	
	\subsection*{Resultados dos testes}

	\begin{itemize}
		\item Funcionamento correto do receiver.
		\par Resultado:
		
		\item Funcionamento correto do transmitter.
		\par Resultado:
		
		\item Verificar funcionamento das portas da entidade UART:
		\begin{table}[H]
			\begin{center}
				\begin{tabular}{|l|l|}
					\hline
					Porta      & Resultado \\ \hline
					A          &           \\ \hline
					notADS     &           \\ \hline
					notBAUDOUT &           \\ \hline
					Din        &           \\ \hline
					Dout       &           \\ \hline
					MR         &           \\ \hline
					RD         &           \\ \hline
					notRXRDY   &           \\ \hline
					SIN        &           \\ \hline
					SOUT       &           \\ \hline
					notTXRDY   &           \\ \hline
					WR         &           \\ \hline
					XIN        &           \\ \hline
					RCLK       &           \\ \hline
				\end{tabular}
			\end{center}
		\end{table}
		
		
	\end{itemize}
	
	
\end{document}
